\documentclass{article}
\usepackage{lmodern}
\usepackage{amsmath}
\usepackage{listings}
\usepackage{fullpage}
\usepackage{parskip}
\usepackage{xcolor}
\lstset{
	basicstyle=\ttfamily,
	columns=fullflexible,
	frame=single,
	breaklines=true,
	postbreak=\mbox{\textcolor{red}{$\hookrightarrow$}\space},
}
\begin{document}
\title{Experiment `p\_4\_5\_analyse\_floats\_array' Results}
\author{\today}
\date{}
\maketitle
\textbf{Experiment outcome: }SUCCESS
\\\textbf{Bad responses: }0
\\\textbf{Responses containing}~\texttt{assume}~\textbf{: }0
\\\textbf{Responses containing}~\texttt{expect}~\textbf{: }0
\\\textbf{Resolution attempts: }1
\\\textbf{Hard fails (resolution): }0
\\\textbf{Soft fails (resolution): }0
\\\textbf{Verification attempts: }1
\section*{Problem Specification}
\textbf{Problem name: }p\_4\_5\_analyse\_floats\_array
\\\textbf{Natural language statement: }Write a method that takes a sequence of floatingpoint numbers and returns: - the average of the values. - the smallest of the values. - the largest of the values. - the range, that is the difference between the smallest and largest.
\\\textbf{Method signature: }p\_4\_5\_analyse\_floats\_array(values: array<real>) returns (average: real, min: real, max: real, range: real)
\subsection*{Ensures}
\begin{itemize}
\item \texttt{average == (sum(values[..]) / values.Length as real)}
\item \texttt{forall v :: v in values[..] ==> min <= v}
\item \texttt{forall v :: v in values[..] ==> max >= v}
\item \texttt{range == max - min}
\item \texttt{exists v :: v in values[..] \&\& min == v}
\item \texttt{exists v :: v in values[..] \&\& max == v}
\end{itemize}
\subsection*{Requires}
\begin{itemize}
\item \texttt{0 < values.Length}
\end{itemize}
\subsection*{Functional Code Given}
\begin{lstlisting}
function sum(values: seq<real>): real
{
  if |values| == 0 then 0.0
  else values[0] + sum(values[1..])
}
\end{lstlisting}
\clearpage
\section*{GenAI interactions}
Below you will find all interactions between the `user' (program) and the `assistant' (OpenAI).
\subsection*{Program $\rightarrow$ GenAI}
\begin{lstlisting}
You are given the following task to perform in Dafny:

Write a method that takes a sequence of floatingpoint numbers and returns: - the average of the values. - the smallest of the values. - the largest of the values. - the range, that is the difference between the smallest and largest.

The signature should be:

p_4_5_analyse_floats_array(values: array<real>) returns (average: real, min: real, max: real, range: real)

The method should respect the following contract:

requires 0 < values.Length, ensures average == (sum(values[..]) / values.Length as real), ensures forall v :: v in values[..] ==> min <= v, ensures forall v :: v in values[..] ==> max >= v, ensures range == max - min, ensures exists v :: v in values[..] && min == v, ensures exists v :: v in values[..] && max == v

The contract uses the following dafny code:

function sum(values: seq<real>): real
{
  if |values| == 0 then 0.0
  else values[0] + sum(values[1..])
}

These function(s) must not be used in your implementation of the method.

Produce and show only the Dafny body of this method, including the curly braces that surround it. Do not show the signature nor contract. You must not use 'assume' nor 'expect' anywhere in your code.
\end{lstlisting}
\subsection*{GenAI $\rightarrow$ Program}
\textbf{System fingerprint: }null
\\\textbf{ID: }chatcmpl-DmfbpGRC3dbqmkIgMtUyDQluGieKN
\\\textbf{Model: }gpt-5.5-2026-04-23
\\\textbf{Created at: }1780492061
\\\textbf{Finish reason: }stop
\\\textbf{Usage: }638 tokens in, and 4461 tokens out
\begin{lstlisting}
{
  var i := values.Length - 1;
  var total := values[i];
  min := values[i];
  max := values[i];

  assert values[i..] == [values[i]];
  assert sum(values[i..]) == values[i];
  assert forall v :: v in values[i..] ==> min <= v;
  assert forall v :: v in values[i..] ==> max >= v;
  assert min in values[i..];
  assert max in values[i..];

  while 0 < i
    invariant 0 <= i < values.Length
    invariant total == sum(values[i..])
    invariant forall v :: v in values[i..] ==> min <= v
    invariant forall v :: v in values[i..] ==> max >= v
    invariant min in values[i..]
    invariant max in values[i..]
    decreases i
  {
    var oldI := i;
    var oldMin := min;
    var oldMax := max;

    i := i - 1;
    var x := values[i];

    assert oldI == i + 1;
    assert values[i..] == [x] + values[oldI..];
    assert values[i..][0] == x;
    assert values[i..][1..] == values[oldI..];

    total := x + total;
    assert total == x + sum(values[oldI..]);
    assert sum(values[i..]) == x + sum(values[i..][1..]);
    assert sum(values[i..]) == total;

    if x < oldMin {
      min := x;
      assert min in values[i..];
      forall v | v in values[i..]
        ensures min <= v
      {
        if v == x {
        } else {
          assert v in values[oldI..];
          assert oldMin <= v;
          assert min < oldMin;
        }
      }
    } else {
      min := oldMin;
      assert min <= x;
      assert min in values[oldI..];
      assert min in values[i..];
      forall v | v in values[i..]
        ensures min <= v
      {
        if v == x {
        } else {
          assert v in values[oldI..];
          assert oldMin <= v;
        }
      }
    }

    if x > oldMax {
      max := x;
      assert max in values[i..];
      forall v | v in values[i..]
        ensures max >= v
      {
        if v == x {
        } else {
          assert v in values[oldI..];
          assert oldMax >= v;
          assert max > oldMax;
        }
      }
    } else {
      max := oldMax;
      assert max >= x;
      assert max in values[oldI..];
      assert max in values[i..];
      forall v | v in values[i..]
        ensures max >= v
      {
        if v == x {
        } else {
          assert v in values[oldI..];
          assert oldMax >= v;
        }
      }
    }
  }

  assert i == 0;
  assert values[i..] == values[..];
  assert total == sum(values[..]);
  assert min in values[..];
  assert max in values[..];

  average := total / (values.Length as real);
  range := max - min;
}
\end{lstlisting}
\section*{Final Program}
\begin{lstlisting}
function sum(values: seq<real>): real
{
  if |values| == 0 then 0.0
  else values[0] + sum(values[1..])
}

method p_4_5_analyse_floats_array(values: array<real>) returns (average: real, min: real, max: real, range: real)
	requires 0 < values.Length
	ensures average == (sum(values[..]) / values.Length as real)
	ensures forall v :: v in values[..] ==> min <= v
	ensures forall v :: v in values[..] ==> max >= v
	ensures range == max - min
	ensures exists v :: v in values[..] && min == v
	ensures exists v :: v in values[..] && max == v
{
  var i := values.Length - 1;
  var total := values[i];
  min := values[i];
  max := values[i];

  assert values[i..] == [values[i]];
  assert sum(values[i..]) == values[i];
  assert forall v :: v in values[i..] ==> min <= v;
  assert forall v :: v in values[i..] ==> max >= v;
  assert min in values[i..];
  assert max in values[i..];

  while 0 < i
    invariant 0 <= i < values.Length
    invariant total == sum(values[i..])
    invariant forall v :: v in values[i..] ==> min <= v
    invariant forall v :: v in values[i..] ==> max >= v
    invariant min in values[i..]
    invariant max in values[i..]
    decreases i
  {
    var oldI := i;
    var oldMin := min;
    var oldMax := max;

    i := i - 1;
    var x := values[i];

    assert oldI == i + 1;
    assert values[i..] == [x] + values[oldI..];
    assert values[i..][0] == x;
    assert values[i..][1..] == values[oldI..];

    total := x + total;
    assert total == x + sum(values[oldI..]);
    assert sum(values[i..]) == x + sum(values[i..][1..]);
    assert sum(values[i..]) == total;

    if x < oldMin {
      min := x;
      assert min in values[i..];
      forall v | v in values[i..]
        ensures min <= v
      {
        if v == x {
        } else {
          assert v in values[oldI..];
          assert oldMin <= v;
          assert min < oldMin;
        }
      }
    } else {
      min := oldMin;
      assert min <= x;
      assert min in values[oldI..];
      assert min in values[i..];
      forall v | v in values[i..]
        ensures min <= v
      {
        if v == x {
        } else {
          assert v in values[oldI..];
          assert oldMin <= v;
        }
      }
    }

    if x > oldMax {
      max := x;
      assert max in values[i..];
      forall v | v in values[i..]
        ensures max >= v
      {
        if v == x {
        } else {
          assert v in values[oldI..];
          assert oldMax >= v;
          assert max > oldMax;
        }
      }
    } else {
      max := oldMax;
      assert max >= x;
      assert max in values[oldI..];
      assert max in values[i..];
      forall v | v in values[i..]
        ensures max >= v
      {
        if v == x {
        } else {
          assert v in values[oldI..];
          assert oldMax >= v;
        }
      }
    }
  }

  assert i == 0;
  assert values[i..] == values[..];
  assert total == sum(values[..]);
  assert min in values[..];
  assert max in values[..];

  average := total / (values.Length as real);
  range := max - min;
}
\end{lstlisting}
\section*{Total Token Usage}
\textbf{Input tokens: }638
\\\textbf{Output tokens: }4461
\\\textbf{Reasoning tokens: }3584
\\\textbf{Sum of `total tokens': }5099
\section*{Experiment Timings}
\textbf{Overall Experiment} started at 1780492061762, ended at 1780492149436, lasting 87674ms (87.67 seconds)
\\\textbf{Iteration \#1} started at 1780492061762, ended at 1780492149436, lasting 87674ms (87.67 seconds)
\end{document}
